% generated by bzoo.report.tables on 2026-09-21; do not edit by hand
\begin{table}[t]
\centering
\caption{Five thresholds for the same family of tests, on one scale. \emph{Nominal} is 1.96. \emph{Calibrated} is the 95th percentile of the measured null for a single test. \emph{Bonferroni} applies the nominal null with $M = 14{,}535$, the number of strategies with a complete 720-month history; the fourth ticker letter position exists only for the most recent 600 months, which is why that count is three quarters of the 19{,}380 and not all of them. The two max-$t$ columns are the 95th percentile of the null maximum with the measured marginals, without and with cross-strategy dependence; the second is correct. \emph{Bonf.\ error} is how much Bonferroni overshoots it. The two alpha rows mix two constructions and should be read as such: \emph{Null SD} and \emph{Calibrated} come from the marginal cross-sectional spread, which Section~\ref{sec:nullfails} shows is not a null, while the max-$t$ columns come from the permutation that imposes $\alpha = 0$ given the exposures. Only the latter is used as a threshold.}
\label{tab:thresholds}
\begin{adjustbox}{max width=\linewidth}
\begin{tabular}{llrrrrrrr}
\toprule
\textbf{Statistic} & \textbf{Wt.} & \textbf{Null SD} & \textbf{Nominal} & \textbf{Calibrated} & \textbf{Bonferroni} & \textbf{Max-$t$ indep.} & \textbf{Max-$t$ joint} & \textbf{Bonf. error} \\
\midrule
Mean return $t$ & EW & 0.889 & 1.96 & 1.77 & 4.64 & 4.73 & 4.39 & 0.057 \\
Five-factor alpha $t$ & EW & 1.150 & 1.96 & 2.27 & 4.64 & 4.35 & 4.12 & 0.127 \\
Mean return $t$ & VW & 1.026 & 1.96 & 2.08 & 4.64 & 4.57 & 4.31 & 0.078 \\
Five-factor alpha $t$ & VW & 1.405 & 1.96 & 3.04 & 4.64 & 4.33 & 4.09 & 0.134 \\
\bottomrule
\end{tabular}
\end{adjustbox}
\end{table}
